\chapter{Methodology, Requirements, and Design Constraints}

The internship produced a designed software artifact for a real organizational
problem: fragile browser evidence collection on volatile streaming websites.
The methodology therefore had to connect problem identification, artifact
design, demonstration, evaluation, and communication. For that reason, this
report adopts Design Science Research Methodology (DSRM) as the sole
methodological frame~\cite{peffers-dsrm,hevner-dsr}.

\section{Methodology Selection}

\subsection{Nature of the Project}

The project was not a pure research survey and not primarily a data-mining or
model-training lifecycle. It produced an artifact: an n8n browser-agent backup
workflow and the OWC browser-agent architecture. The artifact was designed to
address an observed company problem, then demonstrated and evaluated through
real website runs, screenshots, traces, tool calls, provider enrichment, and
regression batches.

\subsection{Why DSRM Fits the Project}

DSRM fits because it provides a clear chain from the organizational problem to
the evaluated artifact. The method starts with the problem and motivation,
defines solution objectives, guides design and development, requires
demonstration, evaluates the artifact, and communicates the result. This maps
directly to the internship work: crawler fragility motivated the project,
browser-agent evidence collection defined the artifact, n8n and OWC formed the
design outputs, real streaming websites demonstrated the approach, and
regression telemetry evaluated it.

\subsection{Why Other Methodologies Were Not Selected}

Agile was not adopted as a methodology in this report. Short correction cycles
were used during implementation, but they are treated as engineering practice
inside the DSRM design-and-evaluation process rather than as a separate
methodology. CRISP-DM was also not selected because the project is not
primarily a data-mining or model-training lifecycle. Batch data and telemetry
are used for evaluation, but the central contribution is an engineering
artifact.

\begin{figure}[!htbp]
\centering
\resizebox{0.94\textwidth}{!}{%
\begin{tikzpicture}[
  card/.style={draw=gray!60, fill=white, rounded corners=6pt,
    text width=3.55cm, minimum height=2.05cm, align=left,
    font=\footnotesize, inner sep=7pt},
  support/.style={card, draw=orange!70!black, fill=orange!8},
  chosen/.style={card, draw=green!55!black, fill=green!8},
  decision/.style={draw=accent!70, fill=blue!5, rounded corners=6pt,
    text width=11.65cm, minimum height=0.9cm, align=center,
    font=\footnotesize\bfseries, inner sep=6pt},
  arr/.style={->, >=Stealth, thick, accent!70, shorten <=3pt, shorten >=3pt}
]
\node[support] (agile) at (-4.65,0) {\centering\textbf{Agile}\par
\vspace{0.25em}\raggedright Short iterations for prompts, tools, browser actions,
and bug fixes.};
\node[support] (crisp) at (0,0) {\centering\textbf{CRISP-DM}\par
\vspace{0.25em}\raggedright Useful for business context, data observation,
evaluation, and operational review.};
\node[chosen] (dsrm) at (4.65,0) {\centering\textbf{DSRM}\par
\vspace{0.25em}\raggedright Main backbone because the contribution is a designed
and evaluated engineering artifact.};

\node[decision] (method) at (0,-2.25)
  {Adopted method: DSRM backbone + Agile iterations + CRISP-DM evaluation lens};

\draw[arr] (agile.south) to[out=-90,in=170] (method.west);
\draw[arr] (crisp.south) -- (method.north);
\draw[arr] (dsrm.south) to[out=-90,in=10] (method.east);
\end{tikzpicture}
}
\caption{Methodology selection logic used for the internship work.}
\label{fig:methodology-selection}
\end{figure}


\section{Adopted Methodology: Design Science Research Methodology}

The report adopts DSRM as the sole methodological frame. The sections below map
the project to the six phases normally used to present DSRM work.

\subsection{Problem Identification and Motivation}

The initial problem was the maintenance burden of existing Puppeteer crawlers
when illegal streaming websites changed. Dynamic pages, popups, hidden iframes,
server switching, tokenized media URLs, and anti-bot behavior made static or
selector-heavy extraction fragile. The motivation was therefore to create a
more traceable browser-agent workflow that could support evidence review.

\subsection{Objectives of the Solution}

The solution objectives were to classify pages, route work to specialist
agents, interact with browser state, capture screenshots, preserve manifests
and media indicators, detect channel or broadcaster context, enrich provider
metadata, and store traces for operator review.

\subsection{Design and Development}

The design phase produced two related artifacts. The n8n workflow validated
browser-agent behavior quickly and served as a backup to the existing crawler
workflow. OWC then formalized the maintainable reference architecture: FastAPI
run control, a Next.js operator console, an orchestrator, specialist agents,
browser tools, persistence, screenshots, provider enrichment, and reviewable
email drafts.

\subsection{Demonstration}

The artifact was demonstrated on real streaming-site cases: landing pages with
repeated event cards, hosting pages with popups and server buttons, embedded
players, iframe handoffs, delayed manifests, tokenized stream URLs, and blocked
or irrelevant pages.

\subsection{Evaluation}

Evaluation used regression batches, screenshots, tool success telemetry,
strict stream success, stream yield, channel/provider coverage, token usage,
cached-token visibility, model spend, and manual review of representative
cases. A run counted as useful only when it produced reviewable evidence or a
defensible traceable failure.

\subsection{Communication}

The final communication artifacts are this PFE report, the diagrams, the
tables, and the evaluation synthesis. They explain the problem, the design
decisions, the implementation, the evaluation evidence, and the remaining
limits.

\begin{figure}[!htbp]
\centering
\resizebox{0.9\textwidth}{!}{%
\begin{tikzpicture}[
  phase/.style={draw=accent!75, fill=blue!5, rounded corners=6pt,
                minimum width=4.0cm, minimum height=1.04cm,
                text centered, font=\small\bfseries, align=center},
  centerbox/.style={draw=orange!70, fill=orange!10, rounded corners=6pt,
                minimum width=3.35cm, minimum height=1.0cm,
                text centered, font=\small\bfseries, align=center},
  arr/.style={->, >=Stealth, thick, accent!72, shorten <=3pt, shorten >=3pt},
  iter/.style={->, >=Stealth, dashed, thick, orange!70!black, shorten <=3pt, shorten >=3pt}
]

\node[phase] (problem) at (0,4.0) {Problem Identification\\and Motivation};
\node[phase] (objectives) at (5.1,2.35) {Objectives of\\the Solution};
\node[phase] (design) at (5.1,-0.75) {Design and\\Development};
\node[phase] (demo) at (0,-2.4) {Demonstration};
\node[phase] (eval) at (-5.1,-0.75) {Evaluation};
\node[phase] (comm) at (-5.1,2.35) {Communication};
\node[centerbox] (loop) at (0,0.75) {Agile-Style\\Short Iterations};

\draw[arr] (problem.east) to[bend left=10] (objectives.north west);
\draw[arr] (objectives.south) -- (design.north);
\draw[arr] (design.south west) to[bend left=10] (demo.east);
\draw[arr] (demo.west) to[bend left=10] (eval.south east);
\draw[arr] (eval.north) -- (comm.south);
\draw[arr] (comm.north east) to[bend left=10] (problem.west);

\draw[iter] (loop.north) -- (problem.south);
\draw[iter] (loop.east) -- (objectives.west);
\draw[iter] (loop.east) -- (design.west);
\draw[iter] (loop.south) -- (demo.north);
\draw[iter] (loop.west) -- (eval.east);
\draw[iter] (loop.west) -- (comm.east);
\end{tikzpicture}
}
\caption{DSRM adopted for the project, with Agile-style short iterations inside each phase.}
\label{fig:dsrm}
\end{figure}


\section{Mapping DSRM to the Internship Work}

\begin{table}[!htbp]
\centering
\begin{tabularx}{\textwidth}{Q{3.1cm} Y Y}
\toprule
\rowcolor{accent!8}
DSRM phase & Internship activity & Output \\
\midrule
Problem identification & Analyze existing crawler fragility, dynamic streaming
sites, and the need for traceable evidence. & Problem statement and motivation. \\
Objectives & Define adaptable browser-agent evidence collection with
screenshots, streams, channel detection, provider enrichment, and traces. &
Functional and non-functional requirements. \\
Design and development & Build the n8n validation/backup workflow and derive
the OWC architecture with agents, tools, persistence, and review surfaces. &
Prototype and reference architecture. \\
Demonstration & Run the workflow on landing, hosting, embedded, popup, iframe,
and tokenized-manifest cases. & Demonstration traces and screenshots. \\
Evaluation & Run regression batches and inspect stream yield, strict success,
tool behavior, provider/channel coverage, and cost telemetry. & Evaluation
tables and failure analysis. \\
Communication & Prepare the PFE report with diagrams, tables, and conclusions.
& Submitted report. \\
\bottomrule
\end{tabularx}
\caption{DSRM mapping to internship activities and outputs.}
\label{tab:dsrm-application}
\end{table}

\section{Project Evolution}

\subsection{Existing Crawler Analysis}

The first phase studied the company crawler context. The crawlers already
performed useful browser automation on known patterns, but their behavior
needed continuous updates when sites changed selectors, player paths, popups,
iframes, server controls, or stream-detection logic.

\subsection{n8n Prototype and Landing-Agent Validation}

The second phase used n8n to validate the browser-agent idea quickly on real
pages. In the company context, the n8n landing agent was deployed as a backup
solution to support the existing crawler workflow. Its role was not to replace
the full crawler system, but to provide a flexible fallback when crawler rules
failed to recover useful landing-page candidates after website changes.

\subsection{Failure Analysis on Dynamic Websites}

The prototype exposed recurring failure modes: popup focus loss, premature
media harvesting, hidden date tabs, repeated event cards, iframe handoffs,
tokenized streams, and pages that mixed landing and hosting signals.

\subsection{Open Web Catcher Reference Architecture}

OWC formalized the approach into a maintainable reference architecture. It
separates the orchestrator, classifier, landing agent, hosting agent, embedded
agent, provider/channel enrichment, email preparation, browser tool layer,
persistence, and operator console.

\subsection{Evaluation and Report Synthesis}

The final phase organized evaluation evidence and report material. The
evaluation remains a snapshot of a volatile target space, so the report avoids
claiming stronger production guarantees than the data supports.

\begin{figure}[!htbp]
\centering
\resizebox{0.98\textwidth}{!}{%
\begin{tikzpicture}[
  marker/.style={circle, fill=accent, draw=accent, minimum size=6pt, inner sep=0pt},
  labelbox/.style={draw=accent!60, fill=blue!6, rounded corners=4pt,
                   minimum width=2.75cm, minimum height=1.0cm,
                   text width=2.45cm, text centered, font=\footnotesize, align=center},
  conn/.style={accent!55, thick, shorten <=2pt, shorten >=2pt}
]
\draw[very thick, accent!60] (0.4,0) -- (15.0,0);
\node[marker] (p1) at (1.0,0) {};
\node[marker] (p2) at (3.8,0) {};
\node[marker] (p3) at (6.6,0) {};
\node[marker] (p4) at (9.4,0) {};
\node[marker] (p5) at (12.2,0) {};
\node[marker] (p6) at (14.6,0) {};

\node[labelbox, fill=gray!10, draw=gray!60] (start) at (1.0,1.35)
  {Evidence\\needs\\analysis};
\node[labelbox, fill=orange!10, draw=orange!60] (proto) at (3.8,-1.35)
  {n8n\\prototype\\runs};
\node[labelbox, fill=orange!10, draw=orange!60] (deploy) at (6.6,1.35)
  {Failure\\analysis\\and fixes};
\node[labelbox, fill=blue!8, draw=accent!70] (python) at (9.4,-1.35)
  {OWC\\platform\\architecture};
\node[labelbox, fill=green!10, draw=green!60] (eval) at (12.2,1.35)
  {Batch\\evaluation\\and telemetry};
\node[labelbox, fill=purple!10, draw=purple!60] (future) at (14.6,-1.35)
  {PFE report\\synthesis};

\draw[conn] (p1.north) -- (start.south);
\draw[conn] (p2.south) -- (proto.north);
\draw[conn] (p3.north) -- (deploy.south);
\draw[conn] (p4.south) -- (python.north);
\draw[conn] (p5.north) -- (eval.south);
\draw[conn] (p6.south) -- (future.north);
\end{tikzpicture}
}
\caption{Internship timeline from evidence needs analysis to evaluated OWC architecture.}
\label{fig:timeline}
\end{figure}


\section{Functional Requirements}

\begin{table}[!htbp]
\centering
\begin{tabularx}{\textwidth}{Q{1.35cm} Y Q{2.1cm}}
\toprule
\rowcolor{accent!8}
ID & Requirement & Priority \\
\midrule
FR-1 & Ingest a target URL and create a traceable investigation run. & Critical \\
FR-2 & Classify the page as landing, hosting, embedded, other, or blocked. & Critical \\
FR-3 & Route work to the correct specialist agent. & Critical \\
FR-4 & Discover landing-page events, channels, match cards, and hosting candidates. & High \\
FR-5 & Operate hosting-page play buttons, server controls, overlays, popups, and iframes. & Critical \\
FR-6 & Inspect embedded-player pages and recover media or manifest evidence. & Critical \\
FR-7 & Detect broadcaster/channel indicators from visible page, DOM, screenshot, stream context, and metadata where possible. & High \\
FR-8 & Capture screenshots at evidence and failure moments. & High \\
FR-9 & Persist \texttt{m3u8}, \texttt{mpd}, \texttt{mp4}, \texttt{webm}, iframe, and media evidence. & Critical \\
FR-10 & Prepare provider and channel-based reviewable email drafts when evidence exists. & Medium \\
FR-11 & Expose evidence bundles and traces for operator review. & High \\
\bottomrule
\end{tabularx}
\caption{Functional requirements for the evidence-collection artifact.}
\label{tab:fr-summary}
\end{table}

\section{Non-Functional Requirements}

\begin{table}[!htbp]
\centering
\begin{tabularx}{\textwidth}{Q{1.35cm} Y Y}
\toprule
\rowcolor{accent!8}
ID & Requirement & Design response \\
\midrule
NFR-1 & Traceability & Store events, tool calls, model calls, screenshots,
outputs, and failure reasons under the same run. \\
NFR-2 & Observability & Keep runtime records that allow engineers to diagnose
failed routing, browser, model, provider, or persistence behavior. \\
NFR-3 & Maintainability & Separate agents, prompts, browser tools,
persistence, configuration, and UI review. \\
NFR-4 & Reliability & Use bounded tools, observe--act--verify discipline,
fail-closed routing, and evidence verification. \\
NFR-5 & Cost control & Track model calls, input tokens, cached input tokens,
output tokens, and estimated spend. \\
NFR-6 & Scalability & Support batch execution and a queue-driven target
deployment direction. \\
NFR-7 & Portability & Keep runtime configuration and browser execution
separable from the API and console. \\
\bottomrule
\end{tabularx}
\caption{Non-functional requirements and design responses.}
\label{tab:nfr-summary}
\end{table}

\section{Design Constraints}

\subsection{Website Volatility}

Layouts, domains, mirror chains, server labels, iframe structures, and media
URLs change frequently. The system must therefore record the live evidence path
instead of trusting memory or old selectors.

\subsection{Browser Runtime Instability}

Browser runs can fail because of timeouts, popups, redirects, blocked
requests, focus changes, or challenge pages. The design must keep those
failures visible.

\subsection{Dynamic Player and Iframe Behavior}

Some pages expose stream evidence only after a click, server switch, iframe
handoff, or wait period. Harvesting too early can create a false failure.

\subsection{Token and Context-Window Pressure}

Browser traces can be long. The agent prompts, tool outputs, memory hints, and
evaluation records must be compact enough to remain usable.

\subsection{Human Review Boundary}

The system prepares evidence and reviewable email drafts. A human reviewer
validates the evidence and decides any downstream action.

\subsection{Ethical and Legal Scope Boundary}

The artifact does not download protected content and does not automatically
send takedown messages. It captures reviewable technical indicators such as
screenshots, manifests, stream candidates, provider context, and traces.

\section{Chapter Summary}

DSRM is the only methodology adopted by this report. Short implementation
corrections happened inside the DSRM process, but they are not presented as a
separate Agile or CRISP-DM methodology. The next chapter introduces the
technical concepts needed to understand the artifact: streaming-page roles,
browser-rendered evidence, media manifests, browser automation, and
traceability.
